\section{证明}
\begin{enumerate}
  \item 证明对象。\zd{哪些事实}需要证据证明。反面讲。
        \begin{kk}{免证事实（不证自明）}
          \begin{enumerate}
            \item 自然规律和定理。不允许推翻。\xjt 可以用相反的证据加以\zd{反驳（内心动摇）\hz 推翻（+确信）}
            \item 众所周知的事实。中国首都是北京。
            \item 根据法律规定推定的事实。署名的人推定为作者。
            \item 根据已知事实和\zd{日常生活经验法则}（区别专业领域知识）所能推定出的另一事实。
            \item 仲裁机构、法院、公证机构\zd{生效文书、已查明}认定的事实。\zd{预决事实}。
          \end{enumerate}
        \end{kk}

        \begin{kk}{自认。重点中的重点}
          \zd{对}于己不利的\zd{事实}陈述/明确表示承认。实践其实很普遍，没有构成\zd{争议焦点}的一般都会承认。
          \begin{enumerate}
            \item \zd{只此三个}场合：庭审中、起诉状答辩状代理词等书面材料、法院组织的调查询问。
            \item 效果：免除对方当事人举证责任。
            \item \zd{认诺}才会当然导致败诉。\zd{对}对方的\zd{诉讼请求}进行承认 \yjt \zd{处分原则}。
            \item 范围：\zd{身份有关的事实}（儿子是隔壁老王的）、程序性事实、公序良俗 \yjt 法院依职权调查。
            \item 方式：明示、默示、委托代理人代为自认（除非 \yi 委托授权书排除代为自认的权利 \er 当庭否认）。
            \item \zd{限制的自认}。
                  \begin{itemize}
                    \item 部分自认。
                    \item 附条件的自认。事实 A + 条件 B 是表达一个完整/两个独立意思，不能断章取义。
                    \item 最高法判断方法：\zd{A、B 是同一个/不同法律关系 \yjt 完整/独立 \yjt 整体/拆分}。
                  \end{itemize}
            \item 共同诉讼人当中的一人自认。\yi 普通 \yjt 谁承认对谁构成自认 \er 必要 \yjt 看其他人态度（\zd{一票否决权}）。
            \item 撤销：法庭辩论终结前，\yi 对方当事人同意 \er 受胁迫 \hz 重大误解。所以想要撤回自认几乎不可能。
            \item 为了达成调解协议 \hz 和解目的所做的\zd{妥协}而认可的事实，都是废话\zd{不算数}。
          \end{enumerate}
        \end{kk}

  \item 证明责任。事实需要证据证明，由\zd{谁来}提供证据证明。本质是结果责任，不利的\zd{风险}负担。
        \begin{itemize}
          \item \zd{得是本案的证明对象}。
          \item 解决一个问题：事实不清真伪不明时，法官如何裁判如何判决，如何认定事实的问题。
          \item 分为：\yi 行为责任：谁提供证据证明；\er \zd{结果责任}：事实不清真伪不明，谁来承担不利后果。
          \item 尽量避免使用证明责任；\zd{是一种拟制，可能与案件的客观事实不符}。
          \item 分配。考试根据原则现场推导。
                \begin{enumerate}
                  \item 一个原则：\zd{谁主张积极事实，谁承担证明责任}。借了 \yjt 积极，没借 \yjt 消极。
                  \item 两个倒置：
                        \begin{enumerate}
                          \item 环境污染（扩大解释，噪声、光污染这种也算）案件因果关系倒置。
                          \item 过错推定案件过错倒置。搁置物悬挂物堆放物等致人损害。
                          \item 过错和免责事由。\zd{活该就是民法中的免责事由}。免责免责，有错才能免，过错是前提。
                        \end{enumerate}
                  \item \zd{做题}：原告下写行为、结果、因果关系、过错，被告下写免责事由。然后一个一个划。穿梭框。
                        % https://q.uiver.app/?q=WzAsNyxbMCwwLCLljp/lkYoiXSxbMiwwLCLooqvlkYoiXSxbMCwxLCLooYzkuLoiXSxbMCwyLCLnu5PmnpwiXSxbMCwzLCLlm6DmnpzlhbPns7siXSxbMCw0LCLml6Dov4fplJkgPyDliKDmjonov4fplJkgOiDnlZnkuIsiXSxbMiwxLCLlhY3otKPkuovnlLEiXSxbMCwxXSxbNCw2LCLnjq/looPmsaHmn5MiLDFdLFs1LDYsIui/h+mUmeaOqOWumiIsMV1d
\[
  \begin{tikzcd}
    \txt{原告} && \txt{被告} \\
    \txt{行为} && \txt{免责事由} \\
    \txt{结果} \\
    \txt{因果关系} \\
    \txt{无过错 ? 删掉过错 : 留下}
    \arrow[from=1-1, to=1-3]
    \arrow["\txt{环境污染}"{description}, from=4-1, to=2-3]
    \arrow["\txt{过错推定}"{description}, from=5-1, to=2-3]
  \end{tikzcd}
\]

                \end{enumerate}
        \end{itemize}

  \item 证明标准。证明到什么\zd{程度}，法院才会相信、采信你主张的事实。法官内心形成的一种确信状态：
        \begin{enumerate}
          \item \zd{一定一定一定要以证明责任为前提}。
          \item 分类。成立的可能性 > 不成立的可能性的。
                \begin{enumerate}
                  \item > 51\% \yjt 较大可能性。程序性事实。
                  \item > 75\% \yjt \zd{高度可能性}（一般的证明标准）。一般的实体性事实。
                  \item > 95\% \yjt 排除合理怀疑。特殊的实体事实：\zd{欺诈、助迫、恶意串通、口头遗属和赠与}。
                \end{enumerate}
        \end{enumerate}
\end{enumerate}
